\chapter{Banach and Hilbert Spaces}

\section*{5.1. The Contraction Mapping Principle}

\begin{definition}\label{gt.def5.1.contraction-mapping}\dcref{ch5:contraction-mapping}\group{chapter-5}
\source{gilbarg-trudinger:page-0086}
\lean{GilbargTrudinger.IsContraction}\leanok
A map $T$ on a normed linear space $\V$ is a contraction if some $\theta<1$
satisfies
\dcref{ch5:5.1}\group{chapter-5}
\begin{equation}\label{gt.eq5.1}
\|Tx-Ty\|\leq\theta\|x-y\| \qquad (x,y\in\V).
\end{equation}
\end{definition}

\begin{theorem}[Contraction mapping principle]\label{gt.thm5.1}\dcref{ch5:5.1}\group{chapter-5}
\source{gilbarg-trudinger:page-0086}
\lean{GilbargTrudinger.contraction_mapping_principle}\leanok
\uses{gt.def5.1.contraction-mapping}
Every contraction $T$ on a Banach space $\B$ has a unique fixed point.
\end{theorem}
\begin{proof}\source{gilbarg-trudinger:page-0086}
\leanok
\sketch For $x_m=T^m x_0$, iteration of \eqref{gt.eq5.1} gives
\[
\|x_n-x_m\|\leq \frac{\theta^m}{1-\theta}\|Tx_0-x_0\|
\qquad(n\geq m).
\]
Completeness yields $x_m\to x$, and continuity gives $Tx=x$. Applying
\eqref{gt.eq5.1} to two fixed points proves uniqueness.
\end{proof}

\section*{5.2. The Method of Continuity}

\begin{definition}\label{gt.def5.2.bounded-linear-operator}\dcref{ch5:bounded-linear-mapping}\group{chapter-5}
\source{gilbarg-trudinger:page-0086,page-0087}
\source{gilbarg-trudinger:page-0087}
\lean{GilbargTrudinger.BoundedLinearOperator, GilbargTrudinger.operatorNormRatios, GilbargTrudinger.operator_norm_eq_sSup_ratio, GilbargTrudinger.bounded_linear_iff_linear_continuous}\leanok
A linear map $T:\V_1\to\V_2$ is bounded when
\dcref{ch5:5.2}\group{chapter-5}
\begin{equation}\label{gt.eq5.2}
\|T\|=\sup_{x\neq0}\frac{\|Tx\|_{\V_2}}{\|x\|_{\V_1}}<\infty.
\end{equation}
For linear maps, boundedness is equivalent to continuity.
\end{definition}

\begin{theorem}[Method of continuity]\label{gt.thm5.2}\dcref{ch5:5.2}\group{chapter-5}
\source{gilbarg-trudinger:page-0087}
\lean{GilbargTrudinger.affineOperator, GilbargTrudinger.method_of_continuity}\leanok
\uses{gt.thm5.1,gt.def5.2.bounded-linear-operator}
Let $\B$ be Banach, let $\V$ be normed, and let
$L_t=(1-t)L_0+tL_1:\B\to\V$ be bounded and linear. If one constant $C$
satisfies
\dcref{ch5:5.3}\group{chapter-5}
\begin{equation}\label{gt.eq5.3}
\|x\|_{\B}\leq C\|L_tx\|_{\V}
\end{equation}
for every $x\in\B$ and $t\in[0,1]$, then $L_0$ is onto if and only if
$L_1$ is onto.
\end{theorem}
\begin{proof}\source{gilbarg-trudinger:page-0087}
\sketch If $L_s$ is onto, \eqref{gt.eq5.3} makes $L_s^{-1}$ well defined
with norm at most $C$. Solving $L_tx=y$ is then the fixed-point problem
\[
x=L_s^{-1}y+(t-s)L_s^{-1}(L_0-L_1)x.
\]
The right side is a contraction whenever
$|t-s|<\{C(\|L_0\|+\|L_1\|)\}^{-1}$. Finitely many such intervals connect
$0$ and $1$.
\end{proof}

\section*{5.3. The Fredholm Alternative}

\begin{definition}\label{gt.def5.3.compact-mapping}\dcref{ch5:compact}\group{chapter-5}
\source{gilbarg-trudinger:page-0087,page-0088}
\lean{GilbargTrudinger.IsCompactMap, GilbargTrudinger.MapsBoundedSequencesToConvergentSubsequences, GilbargTrudinger.compact_map_iff_bounded_sequence, GilbargTrudinger.compact_map_linear_iff_isCompactOperator, GilbargTrudinger.compact_linear_map_continuous}\leanok
A map $T:\V_1\to\V_2$ is compact if it maps bounded sets to relatively
compact sets, equivalently, if every image of a bounded sequence has a
convergent subsequence. A compact linear map is continuous.
\end{definition}

\begin{theorem}[Fredholm alternative]\label{gt.thm5.3}\dcref{ch5:5.3}\group{chapter-5}
\source{gilbarg-trudinger:page-0088,page-0089,page-0090}
\lean{GilbargTrudinger.compact_operator_surjective_sub_smul_id_of_not_hasEigenvalue, GilbargTrudinger.fredholm_alternative}\leanok
\uses{gt.def5.3.compact-mapping,gt.lem5.4}
Let $T:\V\to\V$ be compact and linear. Exactly one of the following holds:
\begin{enumerate}
\item $x-Tx=0$ has a nonzero solution;
\item for every $y\in\V$, the equation $x-Tx=y$ has a unique solution, and
$(I-T)^{-1}$ is bounded.
\end{enumerate}
\end{theorem}
\begin{proof}\source{gilbarg-trudinger:page-0089,page-0090}
\sketch Put $S=I-T$, $\mathcal N=\ker S$, and $\mathcal R=S(\V)$. The
separation argument recorded in Lemma~\ref{gt.lem5.4}, together with
compactness, gives
\dcref{ch5:5.4}\group{chapter-5}\source{gilbarg-trudinger:page-0089}
\begin{equation}\label{gt.eq5.4}
\dist(x,\mathcal N)\leq K\|Sx\|.
\end{equation}
It follows that $\mathcal R$ is closed. If $\mathcal N=\{0\}$ but
$S(\V)\neq\V$, apply the same separation argument to the strictly decreasing
closed ranges $S^j(\V)$; compactness gives a contradiction. Conversely, if
$S$ is onto but has nonzero kernel, apply it to the increasing kernels
$\ker S^j$. Thus injectivity and surjectivity are equivalent, while
\eqref{gt.eq5.4} bounds the inverse.
\end{proof}

\begin{lemma}[Almost orthogonal vectors]\label{gt.lem5.4}\dcref{ch5:5.4}\group{chapter-5}
\source{gilbarg-trudinger:page-0088}
\lean{GilbargTrudinger.almost_orthogonal_vectors}\leanok
If $\M$ is a proper closed subspace of a normed space $\V$, then for every
$\theta<1$ there is $x_\theta\in\V$ such that
$\|x_\theta\|=1$ and $\dist(x_\theta,\M)\geq\theta$.
\end{lemma}
\begin{proof}\source{gilbarg-trudinger:page-0088}
\leanok
\sketch Choose $x\notin\M$, put $d=\dist(x,\M)>0$, and choose
$y_\theta\in\M$ with $\|x-y_\theta\|\leq d/\theta$. Normalizing
$x-y_\theta$ gives the result.
\end{proof}

\begin{definition}\label{gt.def5.3.eigenvalue}\dcref{ch5:eigenvalue}\group{chapter-5}
\source{gilbarg-trudinger:page-0090}
\lean{GilbargTrudinger.IsEigenvalue, GilbargTrudinger.eigenvalueMultiplicity, GilbargTrudinger.compact_operator_bounded_inverse_of_not_eigenvalue}\leanok
\uses{gt.thm5.3}
A scalar $\lambda$ is an eigenvalue of $T$ if $Tx=\lambda x$ for some
$x\neq0$. Its multiplicity is $\dim\ker(\lambda I-T)$. If $T$ is compact,
$\lambda\neq0$, and $\lambda$ is not an eigenvalue, Theorem~\ref{gt.thm5.3}
makes $(\lambda I-T)^{-1}$ a bounded linear map on $\V$.
\end{definition}

\begin{theorem}[Spectrum of a compact operator]\label{gt.thm5.5}\dcref{ch5:5.5}\group{chapter-5}
\source{gilbarg-trudinger:page-0090,page-0091}
\lean{GilbargTrudinger.eigenvalueSet, GilbargTrudinger.spectrum_of_compact_operator}\leanok
\uses{gt.def5.3.compact-mapping,gt.lem5.4,gt.def5.3.eigenvalue}
The eigenvalues of a compact linear endomorphism of a normed space form a
countable set whose only possible accumulation point is $0$. Every nonzero
eigenvalue has finite multiplicity.
\end{theorem}
\begin{proof}\source{gilbarg-trudinger:page-0090,page-0091}
\sketch Otherwise choose linearly independent eigenvectors for eigenvalues
$\lambda_j\to\lambda\neq0$. Lemma~\ref{gt.lem5.4}, applied to their nested
spans, produces unit vectors whose scaled $T$-images stay uniformly
separated, contradicting compactness. The same argument with a repeated
nonzero eigenvalue proves finite multiplicity.
\end{proof}

\section*{5.4. Dual Spaces and Adjoints}

\begin{definition}\label{gt.def5.4.dual-space}\dcref{ch5:dual-space}\group{chapter-5}
\source{gilbarg-trudinger:page-0091}
\lean{GilbargTrudinger.DualSpace, GilbargTrudinger.dual_space_complete, GilbargTrudinger.dual_norm_eq_sSup_ratio}\leanok
The dual $\V^*$ of a normed space $\V$ is the Banach space of bounded linear
functionals with norm
\dcref{ch5:5.5}\group{chapter-5}
\begin{equation}\label{gt.eq5.5}
\|f\|_{\V^*}=\sup_{x\neq0}\frac{|f(x)|}{\|x\|}.
\end{equation}
\end{definition}

\begin{definition}\label{gt.def5.4.reflexive}\dcref{ch5:reflexive}\group{chapter-5}
\source{gilbarg-trudinger:page-0091}
\lean{GilbargTrudinger.canonicalEmbedding, GilbargTrudinger.canonicalEmbedding_apply, GilbargTrudinger.IsReflexiveSpace}\leanok
\uses{gt.def5.4.dual-space}
The canonical map $J:\V\to\V^{**}$ is defined by $Jx(f)=f(x)$. It is a
linear isometry, and $\V$ is reflexive when $J(\V)=\V^{**}$.
\end{definition}

\begin{definition}\label{gt.def5.4.adjoint}\dcref{ch5:adjoint}\group{chapter-5}
\source{gilbarg-trudinger:page-0091}
\lean{GilbargTrudinger.dualAdjoint, GilbargTrudinger.dualAdjoint_apply}\leanok
\uses{gt.def5.2.bounded-linear-operator,gt.def5.4.dual-space}
For a bounded linear map $T:\B_1\to\B_2$, its adjoint
$T^*:\B_2^*\to\B_1^*$ is defined by
\dcref{ch5:5.6}\group{chapter-5}
\begin{equation}\label{gt.eq5.6}
(T^*g)(x)=g(Tx).
\end{equation}
\end{definition}

\section*{5.5. Hilbert Spaces}

\begin{definition}\label{gt.def5.5.inner-product}\dcref{ch5:inner-product}\group{chapter-5}
\source{gilbarg-trudinger:page-0092}
\lean{GilbargTrudinger.RealInnerProductSpace, GilbargTrudinger.inner_product_identities}\leanok
An inner product on a real vector space $\V$ is a symmetric bilinear form
$(\cdot,\cdot)$ with $(x,x)>0$ for $x\neq0$. It induces
$\|x\|=(x,x)^{1/2}$ and satisfies
\dcref{ch5:5.7}\group{chapter-5}
\begin{equation}\label{gt.eq5.7}
|(x,y)|\leq\|x\|\,\|y\|,
\end{equation}
\dcref{ch5:5.8}\group{chapter-5}
\begin{equation}\label{gt.eq5.8}
\|x+y\|\leq\|x\|+\|y\|,
\end{equation}
and
\dcref{ch5:5.9}\group{chapter-5}
\begin{equation}\label{gt.eq5.9}
\|x+y\|^2+\|x-y\|^2=2(\|x\|^2+\|y\|^2).
\end{equation}
\end{definition}

\begin{definition}\label{gt.def5.5.hilbert-space}\dcref{ch5:hilbert-space}\group{chapter-5}
\source{gilbarg-trudinger:page-0092}
\lean{GilbargTrudinger.IsHilbertSpace}\leanok
\uses{gt.def5.5.inner-product}
A Hilbert space is a complete inner-product space.
\end{definition}

\section*{5.6. The Projection Theorem}

\begin{definition}\label{gt.def5.6.orthogonal}\dcref{ch5:orthogonal}\group{chapter-5}
\source{gilbarg-trudinger:page-0093}
\lean{GilbargTrudinger.AreOrthogonal, GilbargTrudinger.orthogonalComplement}\leanok
\uses{gt.def5.5.inner-product}
Vectors $x,y$ are orthogonal when $(x,y)=0$. For $\M\subset\Hspace$, set
$\M^\perp=\{x:(x,y)=0\text{ for every }y\in\M\}$.
\end{definition}

\begin{theorem}[Projection theorem]\label{gt.thm5.6}\dcref{ch5:5.6}\group{chapter-5}
\source{gilbarg-trudinger:page-0093}
\lean{GilbargTrudinger.projection_theorem}\leanok
\uses{gt.def5.5.hilbert-space,gt.def5.6.orthogonal}
If $\M$ is a closed subspace of a Hilbert space $\Hspace$, every
$x\in\Hspace$ has a unique decomposition $x=y+z$ with
$y\in\M$ and $z\in\M^\perp$.
\end{theorem}
\begin{proof}\source{gilbarg-trudinger:page-0093}
\leanok
\sketch The parallelogram identity shows that a minimizing sequence for
$\dist(x,\M)$ is Cauchy, hence converges to some $y\in\M$. Minimality of
$\|x-y\|$ under variations $y+t y'$ implies $(x-y,y')=0$ for all
$y'\in\M$. Orthogonality also proves uniqueness.
\end{proof}

\section*{5.7. The Riesz Representation Theorem}

\begin{theorem}[Riesz representation]\label{gt.thm5.7}\dcref{ch5:5.7}\group{chapter-5}
\source{gilbarg-trudinger:page-0094}
\lean{GilbargTrudinger.riesz_representation}\leanok
For each $F\in\Hspace^*$ there is a unique $f\in\Hspace$ such that
$F(x)=(x,f)$ for all $x\in\Hspace$, and $\|F\|=\|f\|$.
\end{theorem}
\begin{proof}\source{gilbarg-trudinger:page-0094}
\leanok
\uses{gt.thm5.6}
\sketch If $F\neq0$, apply Theorem~\ref{gt.thm5.6} to $\ker F$ and choose
$0\neq z\in(\ker F)^\perp$. Then
$f=F(z)z/\|z\|^2$ represents $F$. Cauchy--Schwarz gives
$\|F\|\leq\|f\|$, and evaluation at $f$ gives the reverse inequality.
\end{proof}

\section*{5.8. The Lax--Milgram Theorem}

\begin{definition}\label{gt.def5.8.bilinear}\dcref{ch5:bounded-coercive-bilinear-form}\group{chapter-5}
\source{gilbarg-trudinger:page-0095}
\lean{GilbargTrudinger.IsBoundedBilinearForm, GilbargTrudinger.IsCoerciveBilinearForm}\leanok
A bilinear form $\mathbf B$ on $\Hilb$ is bounded if some $K$ satisfies
\dcref{ch5:5.10}\group{chapter-5}
\begin{equation}\label{gt.eq5.10}
|\mathbf B(x,y)|\leq K\|x\|\,\|y\|,
\end{equation}
and coercive if some $\nu>0$ satisfies
\dcref{ch5:5.11}\group{chapter-5}
\begin{equation}\label{gt.eq5.11}
\mathbf B(x,x)\geq\nu\|x\|^2.
\end{equation}
\end{definition}

\begin{theorem}[Lax--Milgram]\label{gt.thm5.8}\dcref{ch5:5.8}\group{chapter-5}
\source{gilbarg-trudinger:page-0095}
\lean{GilbargTrudinger.lax_milgram}\leanok
\uses{gt.def5.8.bilinear}
If $\mathbf B$ is bounded and coercive on a Hilbert space $\Hilb$, then for
each $F\in\Hilb^*$ there is a unique $f\in\Hilb$ such that
$\mathbf B(x,f)=F(x)$ for every $x\in\Hilb$.
\end{theorem}
\begin{proof}\source{gilbarg-trudinger:page-0095}
\leanok
\uses{gt.thm5.7}
\sketch Riesz representation defines $T$ by
$\mathbf B(x,f)=(x,Tf)$. Boundedness and coercivity give
$\nu\|f\|\leq\|Tf\|\leq K\|f\|$, so $T$ is injective with closed range.
If the range were proper, a nonzero vector orthogonal to it would contradict
coercivity. Thus $T$ is invertible and the representing vector for $F$ is
$T^{-1}g$.
\end{proof}

\section*{5.9. The Fredholm Alternative in Hilbert Spaces}

\begin{definition}\label{gt.def5.9.adjoint-hilbert}\dcref{ch5:adjoint-in-a-hilbert-space}\group{chapter-5}
\source{gilbarg-trudinger:page-0096}
\lean{GilbargTrudinger.hilbert_adjoint_characterization}\leanok
\uses{gt.def5.5.hilbert-space}
The adjoint of a bounded operator $T:\Hspace\to\Hspace$ is the bounded
operator $T^*$ determined by
\dcref{ch5:5.12}\group{chapter-5}
\begin{equation}\label{gt.eq5.12}
(T^*y,x)=(y,Tx).
\end{equation}
Moreover $\|T^*\|=\|T\|$.
\end{definition}

\begin{lemma}\label{gt.lem5.9}\dcref{ch5:5.9}\group{chapter-5}
\source{gilbarg-trudinger:page-0096}
\lean{GilbargTrudinger.compact_adjoint}\leanok
\uses{gt.def5.3.compact-mapping,gt.def5.9.adjoint-hilbert}
The adjoint of a compact operator on a Hilbert space is compact.
\end{lemma}
\begin{proof}\source{gilbarg-trudinger:page-0096}
\sketch For bounded $x_j$, compactness supplies a Cauchy subsequence of
$TT^*x_j$. The identity
\[
\|T^*(x_j-x_k)\|^2=(x_j-x_k,TT^*(x_j-x_k))
\]
then makes $T^*x_j$ Cauchy.
\end{proof}

\begin{lemma}\label{gt.lem5.10}\dcref{ch5:5.10}\group{chapter-5}
\source{gilbarg-trudinger:page-0096}
\lean{GilbargTrudinger.closure_range_eq_orthogonal_ker_adjoint}\leanok
\uses{gt.def5.9.adjoint-hilbert,gt.thm5.6}
For a bounded operator $T$ on a Hilbert space,
$\overline{\operatorname{ran}T}=(\ker T^*)^\perp$.
\end{lemma}
\begin{proof}\source{gilbarg-trudinger:page-0096}
\leanok
\sketch The adjoint identity gives
$\operatorname{ran}T\subset(\ker T^*)^\perp$. Decompose an arbitrary vector
by the projection theorem relative to $\overline{\operatorname{ran}T}$; its
orthogonal component belongs to $\ker T^*$, which proves the reverse
inclusion after taking orthogonal complements.
\end{proof}

\begin{theorem}[Hilbert-space Fredholm alternative]\label{gt.thm5.11}\dcref{ch5:5.11}\group{chapter-5}
\source{gilbarg-trudinger:page-0096,page-0097}
\lean{GilbargTrudinger.nonzeroEigenvalueSet, GilbargTrudinger.IsRegularValue, GilbargTrudinger.hilbert_space_fredholm_alternative}\leanok
\uses{gt.def5.3.compact-mapping,gt.thm5.5,gt.def5.5.hilbert-space,gt.def5.9.adjoint-hilbert,gt.lem5.9,gt.lem5.10}
Let $T$ be compact on a Hilbert space $\Hspace$. There is a countable
$\Lambda\subset\mathbb R$, with no possible accumulation point except $0$,
such that for $\lambda\neq0$ and $\lambda\notin\Lambda$, both equations
\dcref{ch5:5.13}\group{chapter-5}
\begin{equation}\label{gt.eq5.13}
(\lambda I-T)x=y,\qquad (\lambda I-T^*)x=y
\end{equation}
have unique solutions and bounded inverses. For $\lambda\in\Lambda$, the two
kernels are finite-dimensional and the respective equations are solvable
exactly when $y$ is orthogonal to the kernel of the adjoint operator.
\end{theorem}

\section*{5.10. Weak Compactness}

\begin{definition}\label{gt.def5.10.weak-convergence}\dcref{ch5:weak-convergence}\group{chapter-5}
\source{gilbarg-trudinger:page-0097}
\lean{GilbargTrudinger.WeaklyConverges, GilbargTrudinger.weaklyConverges_iff_inner}\leanok
A sequence $x_j$ in a normed space $\V$ converges weakly to $x$ if
$f(x_j)\to f(x)$ for every $f\in\V^*$. In a Hilbert space this is equivalent
to $(x_j,y)\to(x,y)$ for every $y$.
\end{definition}

\begin{theorem}[Weak sequential compactness]\label{gt.thm5.12}\dcref{ch5:5.12}\group{chapter-5}
\source{gilbarg-trudinger:page-0097}
\lean{GilbargTrudinger.weak_sequential_compactness}\leanok
\uses{gt.def5.5.hilbert-space,gt.thm5.7,gt.def5.10.weak-convergence}
Every bounded sequence in a Hilbert space has a weakly convergent subsequence.
\end{theorem}
\begin{proof}\source{gilbarg-trudinger:page-0097}
\sketch In a separable Hilbert space, diagonalize the scalar sequences
$(x_j,y_m)$ over a countable dense set. Their limits define a bounded linear
functional, represented by some $x$ through Theorem~\ref{gt.thm5.7}; density
then gives $x_{j_k}\rightharpoonup x$. For a general Hilbert space, apply the
separable argument to the closed span of the original sequence and use the
orthogonal decomposition of the ambient space.
\end{proof}

\begin{lemma}\label{gt.lem6.1}\dcref{Lemma 6.1}\group{chapter-6}\source{gilbarg-trudinger:page-0099}
In the equation
\begin{equation}\label{gt.eq6.3}\dcref{(6.3)}\group{chapter-6}\source{gilbarg-trudinger:page-0099}
L_0u=A^{ij}D_{ij}u=f(x),\qquad A^{ij}=A^{ji},
\end{equation}
\end{lemma}
